\section{What is deliberately still wrong}
\label{sec:ch02-still-wrong}

Ask Mosaic for every product, with every product's reviews, and the name of
everyone who wrote one. Call this the catalog-with-reviews query; the rest of
the chapter refers back to it.

\begin{minted}{graphql}
{
  products {
    title
    reviews {
      rating
      author { displayName }
    }
  }
}
\end{minted}

It answers correctly. I timed ten runs of it against a warmed-up service on my
own machine and got a mean of three milliseconds. It also does this:

\begin{minted}{text}
info: Mosaic.ServiceLookups[0]
      Service lookups this request: 146
\end{minted}

\index{N+1 problem}%
One lookup to fetch the product list. Twenty-five more, one per product, to
fetch each product's reviews. A hundred and twenty more, one per review, to
fetch each author. One plus twenty-five plus a hundred and twenty is a hundred
and forty-six round trips through the domain services to answer a single HTTP
request.

Every one of those is a linear scan of a \code{List<T>} that is already in
memory, which is why a hundred and forty-six of them still come in under three
milliseconds. That is exactly why this shape survives into production. Nobody
profiles a service that answers that fast. The lookups stay single-key and
unbatched all the way through review, through staging, and through the first six
months in production. Then somebody puts a database behind the same service, and
a hundred and forty-six list scans become a hundred and forty-six queries against
it, without a line of the service changing.

\index{Mosaic!lookup counter}%
I want to be clear that Mosaic's domain services are written this way on
purpose. Not one of them has a method that takes a list of identifiers. That
omission is the subject of chapter~\ref{ch:data-without-the-n-plus-1}, which
brings EF~Core, explains what a DataLoader does, and brings the same query down
to three lookups. Building the naive version first means that chapter gets to
fix something real rather than demonstrate on a problem invented for the
occasion.

\subsection*{Why the first counter reported 1}

\index{HotChocolate!resolver scopes}%
Counting those lookups took two attempts, and the reason the first one failed is
worth more than the number is.

The obvious design is a scoped service: register a counter per request, inject it
into every domain service, increment it on each lookup, and report the total when
the scope is disposed at the end of the request. That is standard ASP.NET Core,
and in Mosaic it produced this:

\begin{minted}{text}
info: Service lookups this request: 1
info: Service lookups this request: 1
info: Service lookups this request: 1
\end{minted}

A hundred and forty-six times.

HotChocolate resolves a resolver's injected services from a scope of its own,
one per resolver invocation, rather than from the scope that wraps the HTTP
request. Every resolver therefore got a freshly constructed counter, watched it
record exactly one lookup, and disposed it. The total was right in aggregate and
useless in every individual line. Nothing was broken; my mental model of where
the scope boundary sat was simply wrong.

The fix is to stop keeping request state in a scoped service and hang it off the
request instead. A small object goes into \code{HttpContext.Items} in middleware,
the counter increments it through \code{IHttpContextAccessor}, and the middleware
reports the total once the response is finished:

\begin{minted}{csharp}
public static IApplicationBuilder UseServiceCallCounting(this IApplicationBuilder app)
    => app.Use(async (context, next) =>
    {
        var count = new RequestLookupCount();
        context.Items[ServiceCallCounter.ItemKey] = count;

        await next();

        if (count.Value > 0)
        {
            var logger = context.RequestServices
                .GetRequiredService<ILoggerFactory>()
                .CreateLogger("Mosaic.ServiceLookups");

            logger.LogInformation("Service lookups this request: {Count}", count.Value);
        }
    });
\end{minted}

The increment is an \code{Interlocked.Increment}, because resolvers run
concurrently and a hundred and forty-six unsynchronised increments across a
thread pool will quietly lose a few.

The counter itself is throwaway code. The scoping rule it exposed comes back in
chapter~\ref{ch:entity-resolution-done-right}, when DataLoaders have to be shared
across resolvers that are supposedly independent. Resolver lifetime is not
request lifetime, and anything that has to accumulate across a whole GraphQL
request has to live somewhere the request owns.
Chapter~\ref{ch:inside-hotchocolate} takes the execution pipeline apart and shows
where those scopes are actually created.

\subsection*{The rest of the list}

Three other things in this version of Mosaic are wrong on purpose, and naming
them here means the later chapters do not have to pretend to discover them.

\code{reviews} returns a plain list. No pagination, no connection type, no way
to ask for the first ten. Chapter~\ref{ch:schema-design-that-survives-change}
makes it a Relay connection and needs this version to compare against, since a
schema change that nobody can see the before-state of is not much of a lesson.

Errors are exceptions. Pricing throws when it cannot find a price for a
product, and HotChocolate turns that into a GraphQL error with a stack trace
attached in Development. That is fine for seed-data bugs, which is all it can
currently be, and it is not an error model.
Chapter~\ref{ch:schema-design-that-survives-change} builds one.

And nothing is authorised. Any caller can read any customer's orders, given an
identifier, and the seeded identifiers follow a fixed pattern that makes them
easy to guess. Chapter~\ref{ch:identity-and-authorization} fixes this properly
and across service boundaries, which is a harder problem than it looks and not
one to solve twice.

None of that stops Mosaic being a real GraphQL service. It answers queries, it
has a schema you can introspect, it runs in a container, and it is the thing the
next three chapters improve until it is worth distributing.
